\chapter{Energy, geodesics and metric spaces}
\label{chap:q11}
Let $(X,\omega)$ be a connected compact K\"ahler manifold of dimension $n\geq1$. A potential may be unbounded even when its difference from a reference potential is bounded. Non-pluripolar products remove the common poles, but bounded global Stokes cannot be applied directly to the two singular functions. Relative capacity comparison supplies the uniform control needed to pass from bounded truncations. The resulting integration-by-parts identity makes the Monge--Amp\`ere measure the first variation of a scalar energy. Sources are discussed in \cref{sec:history-energy}.

\section{Integration by parts for bounded differences}
\label{sec:singular-bounded-difference-integration}
For a fixed background $\theta$ write $\theta_q=\theta+\ddc q$ and $\mu_q=\langle\theta_q^n\rangle$, indicating another background when needed. A \emph{bounded quasi-psh difference} is a bounded Borel function equal quasi-everywhere to $a=\varphi_1-\varphi_2$, where $\varphi_1,\varphi_2$ are quasi-psh. Thus $\varphi_1\sim\varphi_2$. Values on the common pole set do not enter any integral against a non-pluripolar measure.

For closed positive currents $T_1,\ldots,T_{n-1}$ and $R=\langle T_1\wedge\cdots\wedge T_{n-1}\rangle$, define
\begin{equation}\label{eq:bounded-difference-hessian-product}
 \begin{aligned}
 \ddc a\wedge R
   &=\left\langle(\kappa+\ddc\varphi_1)\wedge T_1\wedge\cdots\wedge T_{n-1}\right\rangle\\
    &\quad-\left\langle(\kappa+\ddc\varphi_2)\wedge T_1\wedge\cdots\wedge T_{n-1}\right\rangle,
 \end{aligned}
\end{equation}
where $\kappa$ is any sufficiently large multiple of a fixed K\"ahler form making both displayed factors positive. This is a finite signed measure. For $n=1$, $R=1$ and the same definition uses the first-degree non-pluripolar product.

This definition is independent of $\kappa$ by positive multilinearity: an added common smooth positive form cancels from the two terms. It is also independent of the presentation of $a$. Indeed, for another presentation $\psi_1-\psi_2$, enlarge $\kappa$ to accommodate all four potentials. The equality
$\varphi_1+\psi_2=\varphi_2+\psi_1$ holds almost everywhere, hence their currents agree. Multilinearity of their positive sums then identifies the two differences in \eqref{eq:bounded-difference-hessian-product}. More generally any smooth closed background making both currents positive gives the same difference, by comparison through a larger common K\"ahler form. We extend the notation linearly to finite sums of mixed products.
\index{quasi-plurisubharmonic function!bounded difference}

\begin{lemma}\label{lem:global-capacity-weighted-measure-convergence}
Suppose positive measures $\nu_j\to\nu$ weakly have uniformly bounded masses, and $\nu_j,\nu$ are uniformly absolutely continuous with respect to $\operatorname{Cap}_\omega$. If uniformly bounded quasi-continuous Borel functions $f_j,f$ satisfy $f_j\to f$ in capacity, then
\[
 \int_Xf_j\,\nu_j\longrightarrow\int_Xf\,\nu.
\]
Uniform absolute continuity means that the supremum of their masses on Borel sets of capacity at most $\delta$ tends to zero as $\delta\downarrow0$.
\end{lemma}
\begin{proof}
For $|f_j|,|f|\leq B$, split the integral of $|f_j-f|$ into the set where it exceeds $\varepsilon$ and its complement. The first part is at most $2B$ times a uniformly small measure mass, by capacity convergence; the second is at most $\varepsilon$ times the common total-mass bound. Thus it tends to zero as $j\to\infty$ and then $\varepsilon\downarrow0$.

For the fixed weight $f$, remove an open set $G$ of arbitrarily small capacity so that $f$ is continuous on $X\setminus G$. The bounded continuous-extension construction in \cref{thm:psh-capacity-quasicontinuity} extends that restriction to a continuous function on $X$ with the same bound. Weak convergence applies to the extension; replacing it by $f$ changes both integrals by at most $2B$ times the measure mass of $G$. Uniform absolute continuity lets this error tend to zero. These two assertions give the result.
\end{proof}

\begin{lemma}\label{lem:singular-monge-ampere-difference-identity}
For a smooth closed real background $\theta$, suppose
$\varphi_1\sim\varphi_2$ and $\psi_1\sim\psi_2$ are $\theta$-psh. Put $a=\varphi_1-\varphi_2$, $b=\psi_1-\psi_2$, and
\[
 S_j=\sum_{k=0}^{n-1}
       \left\langle\theta_{\varphi_j}\wedge\theta_{\psi_1}^{k}
                                      \wedge\theta_{\psi_2}^{n-1-k}\right\rangle,
       \qquad j=1,2.
\]
Then
\begin{equation}\label{eq:singular-monge-ampere-difference-identity}
 \int_Xa(\mu_{\psi_1}-\mu_{\psi_2})
                         =\int_Xb(S_1-S_2).
\end{equation}
No positive-mass assumption on the four original potentials is required.
\end{lemma}
\begin{proof}
First let $\theta=\kappa$ be K\"ahler and suppose all four potentials are $\lambda\kappa$-psh for a common $0<\lambda<1$. Subtract constants to make them negative. Bounded-type mixed mass invariance permits further shifts making
\[
 \varphi_1\leq\varphi_2\leq\varphi_1+B,
                 \qquad\psi_1\leq\psi_2\leq\psi_1+B
\]
with $B\geq1$: the measure differences in \eqref{eq:singular-monge-ampere-difference-identity} have total mass zero, so constant shifts of $a,b$ do not affect the assertion.

Assume first that each pair agrees on an ordinary open set containing its common sufficiently deep sublevel. Truncate all four functions at $-l$. Bounded integration by parts gives \eqref{eq:singular-monge-ampere-difference-identity} for these truncations, by expanding the difference of the two powers. The first measure difference vanishes on the deep open set for the $\psi$ pair, and its weight vanishes on the deep open set for the $\varphi$ pair. On the right, the roles are reversed. Outside these two open sets all four potentials have a common finite lower bound. For large $l$ every remaining weighted term therefore equals the corresponding original non-pluripolar term by locality. The bounded identity is exactly the desired identity.

To remove the agreement assumption take a sequence $\varepsilon\downarrow0$ with $\lambda(1+\varepsilon)<1$ and put
\[
 \varphi_{2,\varepsilon}=\max\{\varphi_1,(1+\varepsilon)\varphi_2\},
 \qquad\psi_{2,\varepsilon}=\max\{\psi_1,(1+\varepsilon)\psi_2\}.
\]
Each function is $\kappa$-psh, stays within $B$ of its first partner and agrees with it on a sufficiently deep ordinary open sublevel. For example $\varphi_1<-B-B/\varepsilon$ implies
$(1+\varepsilon)\varphi_2\leq\varphi_1+B+\varepsilon(\varphi_1+B)<\varphi_1$.
The first case applies to these pairs. As $\varepsilon\downarrow0$ the second partners increase to $\varphi_2,\psi_2$ quasi-everywhere. All their mixed measures converge weakly by increasing continuity.

We need uniform control for the varying bounded weights. Set
\[
 \chi=\tfrac15(\varphi_1+\varphi_2+\psi_1+\psi_2)-B,
 \qquad
 \zeta_\varepsilon=\tfrac15
              (\varphi_1+\varphi_{2,\varepsilon}+\psi_1+\psi_{2,\varepsilon}).
\]
Then $\chi\leq\zeta_\varepsilon\leq\chi+B$, including the limiting $\zeta_0$. Moreover
$\kappa+\ddc\chi\geq(1-4\lambda/5)\kappa$, so $M_\kappa(\chi)>0$. Expanding
\[
 5\kappa_{\zeta_\varepsilon}
    =\kappa_{\varphi_1}+\kappa_{\varphi_{2,\varepsilon}}
                        +\kappa_{\psi_1}+\kappa_{\psi_{2,\varepsilon}}+\kappa
\]
shows that every mixed measure in the identity is bounded by a fixed multiple of $\mu_{\zeta_\varepsilon}$. Strip rescaling bounds this measure on each Borel $E$ by
$B^n\operatorname{Cap}_{\kappa,\chi}(E)$: the potential
$B^{-1}\zeta_\varepsilon+(1-B^{-1})\chi-1$ is a strip test for $\chi$.
The same bounds hold for all limiting measures. The capacity comparison theorem now gives uniform absolute continuity with respect to $\operatorname{Cap}_\omega$.

The differences $\varphi_1-\varphi_{2,\varepsilon}$ and $\psi_1-\psi_{2,\varepsilon}$ are uniformly bounded and quasi-continuous, and converge in capacity to $a,b$. Indeed remove an open set of arbitrarily small capacity containing the poles and making the original four potentials finite and continuous on its compact complement. They are bounded there, so the displayed maxima converge uniformly there. The weighted convergence lemma therefore passes the identity to the limit.

Finally let $\theta$ be arbitrary. For all sufficiently large $s$, $\kappa_s=\theta+s\omega$ is K\"ahler and all four potentials are $\tfrac12\kappa_s$-psh: use $s\omega\geq\theta$. The case already proved gives the identity in background $\kappa_s$. Each side is a polynomial in $s$, by expanding the positive factors $\theta_q+s\omega$ using non-pluripolar multilinearity. Equality for all large $s$ identifies their constant coefficients, which are precisely the two sides of \eqref{eq:singular-monge-ampere-difference-identity} in background $\theta$.
\end{proof}

\Needspace{17\baselineskip}
\begin{theorem}\label{thm:nonpluripolar-bounded-difference-integration-by-parts}
For bounded quasi-psh differences $a,b$ and arbitrary closed positive $(1,1)$-currents $T_1,\ldots,T_{n-1}$,
\begin{equation}\label{eq:nonpluripolar-bounded-difference-integration-by-parts}
 \int_Xa\,\ddc b\wedge\langle T_1\wedge\cdots\wedge T_{n-1}\rangle
       =\int_Xb\,\ddc a\wedge\langle T_1\wedge\cdots\wedge T_{n-1}\rangle.
\end{equation}
The Hessian factors have the meaning of \eqref{eq:bounded-difference-hessian-product}.
\end{theorem}
\begin{proof}
First suppose the presentations of $a,b$ and a potential $\eta$ all use one background $\theta$. Replace $\psi_j$ in the preceding lemma by
$(1-t)\eta+t\psi_j$. Its bounded difference is $tb$. Differentiate the resulting polynomial identity at $t=0$. The derivative of the left side is
$n\int_Xa\,\ddc b\wedge\langle\theta_\eta^{n-1}\rangle$;
the derivative of the right side is
$n\int_Xb\,\ddc a\wedge\langle\theta_\eta^{n-1}\rangle$.
This proves the identity for a repeated current and common background. For $n=1$ it already proves the assertion by choosing a large K\"ahler background and $\eta=0$.

For $n>1$ write $T_j=\theta_j+\ddc\eta_j$. For positive $r_j$ with $\sum_jr_j=1$, put
$\theta_r=\sum_jr_j\theta_j$, $\eta_r=\sum_jr_j\eta_j$.
For all sufficiently large $s$, the presentations of $a,b$ and $\eta_r$ are $\theta_r+s\omega$-psh. Apply the common-background identity to the repeated current
$\sum_jr_jT_j+s\omega$. Presentation independence lets its Hessian factors use one fixed positive background. Expanding the remaining repeated current in $s$ and taking the constant coefficient proves the identity with $\langle(\sum_jr_jT_j)^{n-1}\rangle$. It is a homogeneous polynomial identity in the $r_j$, so validity on the positive simplex implies validity on the entire positive cone by scaling. The coefficient of $r_1\cdots r_{n-1}$, after cancellation of $(n-1)!$, is \eqref{eq:nonpluripolar-bounded-difference-integration-by-parts}. Presentation independence was proved above.
\end{proof}
\index{integration by parts!non-pluripolar products}

\begin{lemma}\label{lem:bounded-difference-negative-hessian-energy}
Suppose $u\sim v$ are $\theta$-psh and $R=\langle T_1\wedge\cdots\wedge T_{n-1}\rangle$. Then
\[
 \int_X(u-v)\,\ddc(u-v)\wedge R\leq0.
\]
\end{lemma}
\begin{proof}
The two measures $\lambda_u=\langle\theta_u\wedge T_1\wedge\cdots\wedge T_{n-1}\rangle$ and $\lambda_v$ have equal total mass by bounded-type mixed mass invariance. Subtract a constant from $u$ to make $u\leq v$; this leaves the integral unchanged. For $t\geq0$ and $\varepsilon>0$, the maximum $q=\max\{u,v-t-\varepsilon\}$ has their same bounded type and mixed mass. Locality on its two strict comparison sets therefore gives
\[
 \lambda_v(\{u<v-t-\varepsilon\})
    \leq\lambda_u(\{u\leq v-t-\varepsilon\})
    \leq\lambda_u(\{u<v-t\}).
\]
Let $\varepsilon\downarrow0$. For every $t\geq0$ this proves
\[
 \lambda_v(\{u<v-t\})\leq\lambda_u(\{u<v-t\}).
\]
Integrate this inequality in $t$. The bounded nonnegative function $v-u$ is the integral of these level indicators, so Tonelli gives
$\int_X(v-u)\lambda_v\leq\int_X(v-u)\lambda_u$.
Rearrangement is the claimed inequality.
\end{proof}

\section{Relative Monge--Amp\`ere energy}
\label{sec:relative-monge-ampere-energy}
Fix a smooth closed real $\theta$ with big class and a $\theta$-model potential $\phi$. Put $m=\int_X\mu_\phi>0$. The relative full-mass class $\mathcal E(X,\theta,\phi)$ was defined in \cref{def:full-nonpluripolar-mass}; all its mixed masses agree with those of $\phi$ by the envelope characterization.

\begin{definition}\label{def:relative-bounded-and-finite-energy-classes}
Write
\[
 \begin{aligned}
 \mathcal E^\infty(X,\theta,\phi)
       &=\{u\in\PSH(X,\theta):u\sim\phi\},\\
 \mathcal E^1(X,\theta,\phi)
       &=\left\{u\in\mathcal E(X,\theta,\phi):
                         \int_X|u-\phi|\,\mu_u<\infty\right\}.
 \end{aligned}
\]
The first class has bounded difference from $\phi$; it need not consist of bounded functions. The second is the \emph{relative finite-energy class}. Polar sets are discarded in the differences. Bounded-type mass invariance gives
$\mathcal E^\infty\subset\mathcal E^1\subset\mathcal E$.
For $\phi=V_\theta$ we omit the reference from these notations.
\end{definition}
\index{finite energy!relative class}\index{relative singularity type!bounded difference}

\begin{definition}\label{def:relative-monge-ampere-energy}
For $u\in\mathcal E^\infty(X,\theta,\phi)$ define
\begin{equation}\label{eq:relative-energy-bounded-type}
 E_\theta^\phi(u)=\frac1{n+1}\sum_{j=0}^n
                  \int_X(u-\phi)\,
                        \langle\theta_u^j\wedge\theta_\phi^{n-j}\rangle.
\end{equation}
All these integrals are finite: the difference is bounded and every mixed measure has total mass $m$. For $\phi=V_\theta$ write $E_\theta$.
\end{definition}
\index{Monge--Ampere energy!relative reference}

\begin{proposition}\label{prop:bounded-type-energy-variation-and-cocycle}
For $u,v\in\mathcal E^\infty(X,\theta,\phi)$, put $u_t=(1-t)u+tv$. Then
\begin{equation}\label{eq:bounded-type-energy-first-variation}
 \frac{\mathrm d}{\mathrm dt}E_\theta^\phi(u_t)
                                   =\int_X(v-u)\,\mu_{u_t}
\end{equation}
with one-sided derivatives at the endpoints, and
\begin{equation}\label{eq:bounded-type-energy-cocycle}
 E_\theta^\phi(v)-E_\theta^\phi(u)
       =\frac1{n+1}\sum_{j=0}^n\int_X(v-u)
                              \langle\theta_v^j\wedge\theta_u^{n-j}\rangle.
\end{equation}
In particular energy is increasing and satisfies
$E_\theta^\phi(u+C)=E_\theta^\phi(u)+Cm$.
\end{proposition}
\begin{proof}
For $w\sim\phi$, positive multilinearity and the elementary beta integral
$\binom nj\int_0^1r^j(1-r)^{n-j}\,\mathrm dr=1/(n+1)$ give
\[
 E_\theta^\phi(w)=\int_0^1\int_X(w-\phi)
                              \mu_{\phi+r(w-\phi)}\,\mathrm dr.
\]
The beta integral follows by repeated integration by parts on the polynomial integrand. All potentials in these mixtures have bounded type $\phi$, so the weights are bounded and all mixed terms finite.

Write $h=v-u$, $q_{r,t}=\phi+r(u_t-\phi)$. The expression for $E_\theta^\phi(u_t)$ is a polynomial in $t$. Differentiating it and applying the integration-by-parts theorem to $u_t-\phi,h$ gives
\[
 \begin{split}
 \frac{\mathrm d}{\mathrm dt}E_\theta^\phi(u_t)
  &=\int_0^1\left[\int_Xh\,\mu_{q_{r,t}}
       +nr\int_X(u_t-\phi)\,\ddc h\wedge
                                      \langle\theta_{q_{r,t}}^{n-1}\rangle\right]\mathrm dr\\
  &=\int_0^1\frac{\mathrm d}{\mathrm dr}
                         \left[r\int_Xh\,\mu_{q_{r,t}}\right]\mathrm dr
   =\int_Xh\,\mu_{u_t}.
 \end{split}
\]
Polynomial differentiation here means differentiation of a finite multilinear expansion; it does not interchange an uncontrolled limit with a singular measure. Integrate from $t=0$ to $1$ and use the same beta integral to obtain \eqref{eq:bounded-type-energy-cocycle}. Positivity proves monotonicity. Every mixed term has mass $m$, so substituting $v=u+C$ gives the shift formula.
\end{proof}

\begin{proposition}\label{prop:bounded-type-energy-concavity}
The energy is concave on $\mathcal E^\infty(X,\theta,\phi)$ and
\begin{equation}\label{eq:bounded-type-energy-tangent-bounds}
 \int_X(v-u)\,\mu_v
       \leq E_\theta^\phi(v)-E_\theta^\phi(u)
       \leq\int_X(v-u)\,\mu_u.
\end{equation}
It is Lipschitz for the uniform norm on differences, with constant $m$. For any other bounded-type reference $\psi\sim\phi$, the same formula defines an energy with
\[
 E_\theta^\psi(u)=E_\theta^\phi(u)-E_\theta^\phi(\psi).
\]
\end{proposition}
\begin{proof}
Differentiate the first variation once more. With $h=v-u$,
\[
 \frac{\mathrm d^2}{\mathrm dt^2}E_\theta^\phi(u_t)
         =n\int_Xh\,\ddc h\wedge
                               \langle\theta_{u_t}^{n-1}\rangle\leq0
\]
by \cref{lem:bounded-difference-negative-hessian-energy}. Thus the derivative decreases; integrating it gives the two tangent bounds. The first variation and $M(u_t)=m$ give
$|E_\theta^\phi(v)-E_\theta^\phi(u)|\leq m\|v-u\|_\infty$, where the norm ignores the polar set. Finally apply the cocycle formula with $u=\psi$ to identify the energy based at $\psi$.
\end{proof}

\begin{definition}\label{def:extended-relative-energy}
For $u\in\PSH(X,\theta)$ with $u\preceq\phi$, extend energy by
\[
 E_\theta^\phi(u)=\inf\{E_\theta^\phi(w):w\in\mathcal E^\infty(X,\theta,\phi),\ u\leq w\}.
\]
The value lies in $[-\infty,\infty)$, since $\phi+C$ is a majorant for sufficiently large $C$. Monotonicity shows that this agrees with \eqref{eq:relative-energy-bounded-type} on bounded type.
\end{definition}

\Needspace{16\baselineskip}
\begin{proposition}\label{prop:relative-energy-canonical-truncations}
The extended energy is increasing, has the same constant-shift formula, and satisfies
\[
 E_\theta^\phi(u)=\lim_{k\to\infty}
                       E_\theta^\phi(\max\{u,\phi-k\}).
\]
In particular the limit exists, allowing $-\infty$.
\end{proposition}
\begin{proof}
Monotonicity follows by inclusion of the majorant families; translating a majorant by a constant proves the shift identity. Write $u_k=\max\{u,\phi-k\}$. They are decreasing bounded-type majorants of $u$, so their energies decrease and bound $E_\theta^\phi(u)$ above. For any other bounded-type majorant $w$, choose $C$ with $w\geq\phi-C$. For $k\geq C$, both $u$ and $\phi-k$ are at most $w$, so $u_k\leq w$ and $E_\theta^\phi(u_k)\leq E_\theta^\phi(w)$. Taking the limit and then the infimum over $w$ gives the reverse inequality.
\end{proof}

\begin{proposition}\label{prop:relative-energy-decreasing-continuity}
For a decreasing net $u_i\in\PSH(X,\theta)$ with $u_i\preceq\phi$ and potential limit $u=\inf_i u_i$,
\[
 E_\theta^\phi(u_i)\downarrow E_\theta^\phi(u).
\]
Consequently the extended energy is concave. Its finite locus is convex.
\end{proposition}
\begin{proof}
First consider a sequence. Monotonicity gives a limit $L\geq E_\theta^\phi(u)$. For any bounded-type majorant $w\geq u$, put $v_i=\max\{u_i,w\}$. Then $v_i\downarrow w$ and all $v_i$ have a common bounded difference from $w$: use $u_i\leq u_1\leq\phi+C$ and $w\sim\phi$. The tangent bound gives
\[
 0\leq E_\theta^\phi(v_i)-E_\theta^\phi(w)
                         \leq\int_X(v_i-w)\,\mu_w\longrightarrow0
\]
by bounded convergence. Since $E_\theta^\phi(u_i)\leq E_\theta^\phi(v_i)$, taking the infimum over $w$ gives $L\leq E_\theta^\phi(u)$.

For a net, global decreasing compactness gives $L^1$ convergence to $u$. Choose increasing indices $i_k$ with $\|u_{i_k}-u\|_{L^1}<1/k$. Their decreasing pointwise limit is a potential equal to $u$ almost everywhere, hence everywhere. The sequence case gives $E_\theta^\phi(u_{i_k})\to E_\theta^\phi(u)$. The infimum of the full energy net is no greater than this selected limit, proving equality; no countable cofinality is asserted.

For concavity use the bounded-type mixtures of the canonical truncations of $u,v$. They decrease to $(1-t)u+tv$ for $0<t<1$. Bounded-type concavity and decreasing continuity give
\[
 E_\theta^\phi((1-t)u+tv)
                    \geq(1-t)E_\theta^\phi(u)+tE_\theta^\phi(v).
\]
This also proves convexity of the finite locus. The assertion is automatic if either energy on the right is $-\infty$.
\end{proof}

\Needspace{21\baselineskip}
\begin{theorem}\label{thm:relative-finite-energy-characterization}
For $u\preceq\phi$,
\[
 E_\theta^\phi(u)>-\infty
                    \quad\Longleftrightarrow\quad
                          u\in\mathcal E^1(X,\theta,\phi).
\]
In particular finite extended energy forces relative full mass. For $u\leq\phi$ in this class,
\begin{equation}\label{eq:relative-energy-integrability-bounds}
 \frac1{n+1}\int_X(\phi-u)\,\mu_u
       \leq -E_\theta^\phi(u)
       \leq\int_X(\phi-u)\,\mu_u.
\end{equation}
\end{theorem}
\begin{proof}
Subtract a constant to make $u\leq\phi$; both conditions are invariant under this operation. Put $u_k=\max\{u,\phi-k\}$ and $A_k=\{u>\phi-k\}$. These potentials have mass $m$. The energy definition and nonpositive weights give
\[
 \int_X(\phi-u_k)\mu_{u_k}\leq(n+1)(-E_\theta^\phi(u_k)).
\]
If $E_\theta^\phi(u)>-\infty$, the right side is bounded by a fixed $C$. On $X\setminus A_k$ the weight is $k$ outside a polar set; hence
\[
 m-\mu_u(A_k)=\mu_{u_k}(X\setminus A_k)\leq C/k.
\]
Locality gives the equality. The sets $A_k$ increase to the common finite locus, so $M(u)=m$. The same estimate and locality give
$\int_{A_k}(\phi-u)\mu_u\leq C$; monotone convergence proves finite self-integrability and the first inequality in \eqref{eq:relative-energy-integrability-bounds}.

Conversely suppose $u\in\mathcal E^1$ and let $I=\int_X(\phi-u)\mu_u<\infty$. Equality of masses and locality yield
$\mu_{u_k}(X\setminus A_k)=\mu_u(X\setminus A_k)$. Thus
\[
 \int_X(\phi-u_k)\mu_{u_k}
      =\int_X\min\{\phi-u,k\}\,\mu_u\leq I.
\]
The lower tangent bound for $u_k,\phi$ gives $E_\theta^\phi(u_k)\geq-I$. Passing to the canonical energy limit proves finite energy and the second inequality. The first inequality follows by the same finite-level estimate already proved.
\end{proof}
\index{finite energy!characterization by extended energy}

\begin{example}\label{ex:curve-dirichlet-and-loglog-energy}
On a compact K\"ahler curve with reference zero, a smooth $\omega$-psh potential has
\[
 E_\omega(u)=\int_Xu\,\omega-\tfrac12
                         \int_X\mathrm du\wedge\mathrm d^{\mathrm c}u.
\]
Indeed the two terms in the energy definition are $\int u\omega$ and $\int u(\omega+\ddc u)$, and bounded integration by parts gives the formula. Thus in dimension one the relative energy contains the ordinary Dirichlet energy.

The unbounded full-mass potential on $\mathbb P^1$ in \cref{ex:loglog-full-mass-envelope} also lies in $\mathcal E^1(\mathbb P^1,\omega,0)$. Near its only pole, write $r=\log|z|^2$. There $F(r)=-\log(-r)$, $F''(r)=r^{-2}$, and the radial pushforward of $\mu_u=\ddc F(\log|z|^2)$ is $F''(r)\,\mathrm dr$, by the normalization in \cref{thm:euclidean-jensen-lelong}. The smooth correction in $u$ is bounded there, so for $A>1$ its self-integrability reduces to
\[
 \int_A^\infty\frac{\log t}{t^2}\,\mathrm dt
                         =\frac{\log A+1}{A}<\infty.
\]
Away from the pole $u$ is bounded and the measure finite. Hence this is a finite-energy potential with unbounded difference from its model reference zero.
\end{example}
\index{finite energy!log-log example}\index{Dirichlet energy!Monge--Ampere energy on curves}

\Needspace{14\baselineskip}
\begin{corollary}\label{cor:relative-energy-classes-convex}
The classes $\mathcal E^\infty$, $\mathcal E^1$ and $\mathcal E$ relative to $\phi$ are convex. If $u\in\mathcal E^1$ and $u\leq v\preceq\phi$, then $v\in\mathcal E^1$.
\end{corollary}
\begin{proof}
Bounded type is preserved by convex combinations. Finite energy is preserved by concavity, and the characterization proves the assertion for $\mathcal E^1$. For $u,v\in\mathcal E$, their $P_\theta$ representatives equal $\phi$. Hence every mixed mass in the expansion of $\mu_{(1-t)u+tv}$ is $m$, by mixed envelope mass preservation. The total mass is $m$, proving the assertion for $\mathcal E$. Monotonicity of extended energy proves the last claim.
\end{proof}

\begin{lemma}\label{lem:relative-cross-energy-estimate}
If $u,v,\gamma\in\mathcal E(X,\theta,\phi)$ and $\gamma\geq\max\{u,v\}$, then, allowing infinite values,
\[
 \int_X(\gamma-u)\mu_v
       \leq2\int_X(\gamma-u)\mu_u
                          +2\int_X(\gamma-v)\mu_v.
\]
Consequently for $u,v,w_1,\ldots,w_n\in\mathcal E^1(X,\theta,\phi)$,
\[
 \int_X|u-v|\,\langle\theta_{w_1}\wedge\cdots\wedge\theta_{w_n}\rangle<\infty.
\]
\end{lemma}
\begin{proof}
For $t\geq0$,
\[
 \{\gamma>u+2t\}\subset
                     \{\gamma>v+t\}\cup\{v>u+t\}.
\]
Layer integration and the full-mass comparison on the second set give
\[
 \begin{split}
 \int_X(\gamma-u)\mu_v
  &\leq2\int_0^\infty\mu_v(\{\gamma>v+t\})\,\mathrm dt
          +2\int_0^\infty\mu_u(\{v>u+t\})\,\mathrm dt\\
  &\leq2\int_X(\gamma-v)\mu_v+2\int_X(\gamma-u)\mu_u.
 \end{split}
\]
All weights are nonnegative, so Tonelli applies without an integrability assumption.

For the second claim shift $u,v$ below $\phi$ and put $w=n^{-1}\sum_jw_j\in\mathcal E^1$. Positive multilinearity bounds the mixed measure by $n^n\mu_w$. Choose $C$ with $w\leq\phi+C$ and apply the first claim to $u,w,\phi+C$. Its right side is finite by self-integrability of $u,w$. Hence $\phi-u$ is integrable for the mixed measure. Apply the same argument to $v$ and use $|u-v|\leq(\phi-u)+(\phi-v)$ after the shifts; restoring constants adds only a finite multiple of $m$.
\end{proof}

\begin{proposition}\label{prop:relative-finite-energy-cocycle}
For $u,v\in\mathcal E^1(X,\theta,\phi)$, the cocycle formula \eqref{eq:bounded-type-energy-cocycle} and tangent bounds \eqref{eq:bounded-type-energy-tangent-bounds} remain valid. In particular \eqref{eq:relative-energy-bounded-type} computes extended energy throughout $\mathcal E^1$.
\end{proposition}
\begin{proof}
Shift $u,v$ below $\phi$ and put $u_k=\max\{u,\phi-k\}$, $v_k=\max\{v,\phi-k\}$. Fix $0\leq j\leq n$ and write
$\lambda_k=\langle\theta_{u_k}^j\wedge\theta_{v_k}^{n-j}\rangle$ and $\lambda=\langle\theta_u^j\wedge\theta_v^{n-j}\rangle$.
On $A_k=\{u>\phi-k\}\cap\{v>\phi-k\}$ they agree and $u_k-v_k=u-v$. The preceding lemma makes $|u-v|$ integrable for $\lambda$, so its integral on $X\setminus A_k$ tends to zero. It remains to bound the truncated tail.

For either $z=u,v$ and $D_k(z)=\{z\leq\phi-k\}$ we claim
\begin{equation}\label{eq:finite-energy-mixed-truncation-tail}
 k\lambda_k(D_k(z))\leq
                         2^n k\mu_z(\{z\leq\phi-k/2\})\longrightarrow0.
\end{equation}
For $z=u$, put $q=(v_k+\phi-k)/2$. On $D_k(u)$,
$u_k\leq q\leq\phi-k/2$. Positive multilinearity bounds
$\lambda_k\leq2^{n-j}\langle\theta_{u_k}^j\wedge\theta_q^{n-j}\rangle$.
The non-strict mixed full-mass comparison for $q,u_k$, with $\theta_{u_k}^j$ as the fixed factors, therefore gives
\[
 \lambda_k(D_k(u))
   \leq2^{n-j}\mu_{u_k}(\{u_k\leq q\})
   \leq2^{n-j}\mu_{u_k}(\{u\leq\phi-k/2\}).
\]
For $j=n$ the same bound follows directly, with no comparison factors. Total-mass equality and locality on $\{u>\phi-k/2\}$ identify the last mass with $\mu_u(\{u\leq\phi-k/2\})$. The argument with $u,v$ interchanged proves the other case. Finally
$k\mu_z(\{z\leq\phi-k/2\})\leq2\int_{\{z\leq\phi-k/2\}}(\phi-z)\mu_z\to0$
by self-integrability, proving \eqref{eq:finite-energy-mixed-truncation-tail}.

Both $u_k-\phi$ and $v_k-\phi$ lie in $[-k,0]$. Thus $|u_k-v_k|\leq k$, and the tail bound proves
\[
 \int_X(u_k-v_k)\lambda_k\longrightarrow\int_X(u-v)\lambda.
\]
Pass to the limit in the bounded-type cocycle, using decreasing energy continuity. The same argument with $j=n$ and $j=0$ passes its tangent bounds. Restore the constants using equality of mixed masses. Taking $v=\phi$ gives the energy formula.
\end{proof}

\Needspace{23\baselineskip}
\section{Continuity and rooftop stability}
\label{sec:relative-energy-stability}
The reference $\phi$ remains fixed. Increasing limits require control of the truncation error, while $L^1$ limits are naturally handled by decreasing upper envelopes. Both arguments will also be useful for the metric topology.

\begin{lemma}\label{lem:ordered-relative-self-integral-estimate}
If $u\leq v\leq\gamma$ belong to $\mathcal E(X,\theta,\phi)$, then, allowing infinite values,
\[
 \int_X(\gamma-v)\,\mu_v
                 \leq2^{n+1}\int_X(\gamma-u)\,\mu_u.
\]
\end{lemma}
\begin{proof}
Put $q=(\gamma+v)/2$. This has relative full mass by convexity. For $t>0$, outside the common polar exceptional set,
\[
 \{\gamma>v+2t\}=\{q>v+t\}
       \subseteq\{q>u+t\}\subseteq\{\gamma>u+t\}.
\]
Positive multilinearity gives $\mu_v\leq2^n\mu_q$. Apply relative comparison to $u$ and $q-t$ on the middle set, then integrate its mass bound over $t\in(0,\infty)$. The layer formula for the left side has the additional factor $2$. Tonelli's theorem proves the assertion without any preliminary integrability assumption.
\end{proof}

\begin{proposition}\label{prop:relative-energy-increasing-continuity}
Let $(u_i)_{i\in I}$ be an increasing net in $\mathcal E^1(X,\theta,\phi)$ with a common upper bound on $X$. Then its canonical envelope $u=(\sup_i u_i)^*$ belongs to $\mathcal E^1(X,\theta,\phi)$ and
\[
 E_\theta^\phi(u_i)\uparrow E_\theta^\phi(u).
\]
\end{proposition}
\begin{proof}
Model normalization, \cref{lem:model-normalized-supremum}, gives a common bound $u_i\leq\phi+C$. The same holds for $u$. It dominates any fixed member of the family, so monotonicity and the finite-energy characterization put it in $\mathcal E^1$.

First consider an increasing sequence, and subtract a common constant so that $u_j\leq u\leq\phi$. If all functions lie between $\phi-C$ and $\phi$, with $C\geq1$, then
$\mu_{u_j}\leq C^n\operatorname{Cap}_{\theta,\phi}$ as set functions, by the strip rescaling used earlier. The same holds for $\mu_u$. Capacity comparison gives uniform absolute continuity with respect to $\operatorname{Cap}_\omega$. Increasing continuity gives weak convergence of the measures. The bounded quasi-continuous differences $u-u_j$ tend to zero in capacity: use local extended-capacity convergence and remove an open set of arbitrarily small capacity containing the poles and making $\phi$ finite and continuous on the remaining compact set. There the common lower bound $\phi-C$ is finite and bounded, so all the differences are ordinary bounded differences. The weighted convergence lemma and the tangent bound yield
\[
 0\leq E_\theta^\phi(u)-E_\theta^\phi(u_j)
                  \leq\int_X(u-u_j)\,\mu_{u_j}\longrightarrow0.
\]

For the general sequence fix its first member $u_1$, and set $u_{j,k}=\max\{u_j,\phi-k\}$ and $u_k=\max\{u,\phi-k\}$. For fixed $k$ the preceding case gives $E_\theta^\phi(u_{j,k})\to E_\theta^\phi(u_k)$. The finite-energy tangent bound gives
\[
 0\leq E_\theta^\phi(u_{j,k})-E_\theta^\phi(u_j)
       \leq\int_k^\infty\mu_{u_j}\{u_j<\phi-t\}\,\mathrm dt.
\]
For $q_t=(u_j+\phi-t)/2$, the set on the right is contained in $\{u_1<q_t\}$. Multilinearity and relative comparison show
\[
 \mu_{u_j}\{u_j<\phi-t\}
 \leq2^n\mu_{q_t}\{u_1<q_t\}
 \leq2^n\mu_{u_1}\{u_1<q_t\}
 \leq2^n\mu_{u_1}\{u_1<\phi-t/2\}.
\]
Consequently
\begin{equation}\label{eq:relative-energy-uniform-truncation-error}
 \begin{aligned}
 0&\leq E_\theta^\phi(u_{j,k})-E_\theta^\phi(u_j)\\
  &\leq2^{n+1}\int_X(\phi-u_1-k/2)_+\,\mu_{u_1}\longrightarrow0
 \end{aligned}
\end{equation}
uniformly in $j$. Self-integrability justifies the last limit. Let $j\to\infty$ and then $k\to\infty$, using canonical energy truncation, to obtain the sequence assertion.

For a net, choose a countable subfamily with the same canonical upper envelope by \cref{thm:psh-upper-envelopes}. Directedness supplies increasing indices $i_k$ dominating its first $k$ indices and one fixed initial index. The selected sequence has canonical envelope $u$ and tends to it quasi-everywhere by \cref{cor:psh-pluripolar-discrepancy}. Its energies tend to $E_\theta^\phi(u)$. Every one of those energies is at most the supremum of the full net, while monotonicity bounds that supremum by $E_\theta^\phi(u)$. This proves equality; no countable cofinality is asserted.
\end{proof}
\index{Monge--Amp\`ere energy!increasing continuity}

\begin{proposition}\label{prop:relative-energy-lone-upper-semicontinuity}
The set $\{u\in\PSH(X,\theta):u\preceq\phi\}$ is closed in the $L^1$ topology of potentials, and $E_\theta^\phi$ is upper semicontinuous there, with values in $[-\infty,\infty)$. Thus if $u_j\to u$ in $L^1$, then
\[
 E_\theta^\phi(u)\geq\limsup_j E_\theta^\phi(u_j).
\]
\end{proposition}
\begin{proof}
Continuity of the supremum in \cref{prop:normalized-global-psh-compactness} and model normalization give a common bound $u_j\leq\phi+C$. Almost-everywhere subsequence convergence, followed by uniqueness of the canonical representative, gives $u\leq\phi+C$.

Set $v_j=(\sup_{k\geq j}u_k)^*$. These are potentials with $v_j\leq\phi+C$, and they decrease to $u$. To check the last assertion directly, almost-everywhere subsequence convergence gives $u\leq v_j$. For a finite value $u(x)$ and $\varepsilon>0$, choose a neighborhood where $u<u(x)+\varepsilon$. The continuous-majorant version of Hartogs' lemma, after adding a smooth local background potential, bounds all sufficiently late $u_k$ by $u(x)+\varepsilon$ on a smaller neighborhood. Their regularized tail envelope obeys the same bound at $x$. At a pole use any finite upper bound in place of $u(x)+\varepsilon$. Thus $\inf_jv_j=u$ pointwise. Decreasing energy continuity now gives
\[
 \limsup_j E_\theta^\phi(u_j)
       \leq\lim_jE_\theta^\phi(v_j)=E_\theta^\phi(u).
\]
Sequential upper semicontinuity is upper semicontinuity because $L^1$ is a metric topology.
\end{proof}

\begin{corollary}\label{cor:relative-energy-compact-sublevels}
For $B\geq0$, the set
\[
 \{u\in\PSH(X,\theta):u\leq\phi,\quad E_\theta^\phi(u)\geq-B\}
\]
is compact in $L^1$ and lies in $\mathcal E^1(X,\theta,\phi)$.
\end{corollary}
\begin{proof}
Model normalization gives $u\leq\phi+\sup_Xu$, hence
$E_\theta^\phi(u)\leq m\sup_Xu$. Every member therefore satisfies $-B/m\leq\sup_Xu\leq0$. Global compactness supplies subsequential potential limits. The order bound persists in $L^1$, and upper semicontinuity preserves the lower energy bound. The finite-energy characterization gives membership in $\mathcal E^1$. Closedness inside the compact supremum interval proves compactness.
\end{proof}
\index{finite energy!compact energy sublevels}

\Needspace{19\baselineskip}
\begin{lemma}\label{lem:relative-self-integral-lone-stability}
Suppose $\gamma,u_j\in\mathcal E(X,\theta,\phi)$, $u_j\leq\gamma$, and $u_j\to u\in\PSH(X,\theta)$ in $L^1$. If
\[
 \int_X(\gamma-u_j)\,\mu_{u_j}\leq A
\]
for all $j$, then $u\in\mathcal E(X,\theta,\phi)$ and
\[
 \int_X(\gamma-u)\,\mu_u\leq(n+1)A.
\]
The reference $\gamma$ need not have finite energy relative to $\phi$.
\end{lemma}
\begin{proof}
The order bound passes to $u$. On potentials $w\sim\gamma$ define the auxiliary bounded-type energy
\[
 F_\gamma(w)=\frac1{n+1}\sum_{a=0}^n
       \int_X(w-\gamma)
           \left\langle\theta_w^a\wedge\theta_\gamma^{n-a}\right\rangle.
\]
The proofs of \cref{prop:bounded-type-energy-variation-and-cocycle,prop:bounded-type-energy-concavity} apply to this formula: their polynomial calculation uses only integration by parts and equality of bounded-type mixed masses, not the model property of the reference. In particular $F_\gamma(\gamma)=0$, and
\[
 F_\gamma(w)\geq-\int_X(\gamma-w)\,\mu_w
                       \qquad(w\leq\gamma).
\]
For potentials between $\gamma-k$ and $\gamma$, this energy is $L^1$ upper semicontinuous. Indeed regularized tail suprema decrease to the $L^1$ limit by the Hartogs argument above, stay in the same interval, and monotonicity bounds the original energies by their energies. For a decreasing sequence $z_l\downarrow z$ in that interval, the tangent bound gives
\[
 0\leq F_\gamma(z_l)-F_\gamma(z)
                       \leq\int_X(z_l-z)\,\mu_z\longrightarrow0
\]
by dominated convergence. This proves the asserted upper semicontinuity on the interval.

Put $u_{j,k}=\max\{u_j,\gamma-k\}$ and $u_k=\max\{u,\gamma-k\}$. The former have mass $m$ by bounded-type invariance. Since $u_j$ also has mass $m$, locality on $\{u_j>\gamma-k\}$ and subtraction from the equal total masses give
\[
 \int_X(\gamma-u_{j,k})\,\mu_{u_{j,k}}
       =\int_X\min\{\gamma-u_j,k\}\,\mu_{u_j}\leq A.
\]
They tend in $L^1$ to $u_k$, since taking a maximum with the fixed almost-everywhere finite $\gamma-k$ is Lipschitz for $L^1$ differences. Thus $F_\gamma(u_k)\geq-A$. All terms of its definition are nonpositive, so
\[
 \int_X(\gamma-u_k)\,\mu_{u_k}\leq(n+1)A.
\]
On the discarded set $\{u\leq\gamma-k\}$ the weight is $k$. Locality on its complement therefore gives
\[
 0\leq m-\mu_u\{u>\gamma-k\}\leq(n+1)A/k.
\]
Let $k\to\infty$. All poles have zero non-pluripolar mass, so $\mu_u(X)=m$. The same locality and monotone convergence pass the integral bound to $\int_X(\gamma-u)\,\mu_u$. Finally $u\leq\gamma\preceq\phi$, which proves its relative full-mass membership.
\end{proof}

\begin{proposition}\label{prop:relative-energy-family-supremum}
A nonempty family in $\mathcal E(X,\theta,\phi)$, $\mathcal E^1(X,\theta,\phi)$ or $\mathcal E^\infty(X,\theta,\phi)$ with a common upper bound has its canonical supremum in the same class.
\end{proposition}
\begin{proof}
Model normalization bounds every member by $\phi+C$ with a common $C$. The canonical supremum obeys that bound and dominates a chosen member. Mass monotonicity proves the full-mass assertion; energy monotonicity and the characterization prove the finite-energy assertion. In the bounded-type case the chosen member supplies a lower bound $\phi-C'$, proving the last assertion.
\end{proof}

\Needspace{14\baselineskip}
\begin{proposition}\label{prop:relative-finite-energy-rooftop-closure}
For $u,v\in\mathcal E^1(X,\theta,\phi)$, their rooftop $w=P_\theta(u,v)$ belongs to $\mathcal E^1(X,\theta,\phi)$. If $u,v\leq\phi$, then
\[
 \int_X(\phi-w)\,\mu_w
     \leq\int_X(\phi-u)\,\mu_u+\int_X(\phi-v)\,\mu_v.
\]
\end{proposition}
\begin{proof}
\Cref{prop:positive-mass-rooftop-existence,thm:rooftop-mass-diamond} prove that $w$ exists and has mass $m$: $\max\{u,v\}$ has mass $m$, so the existence margin is $2m>m$, and the diamond and mass monotonicity give $M(w)=m$. Subtract a common constant so that $u,v\leq\phi$. The contact inequality gives
\[
 \mu_w\leq\mathbf1_{\{w=u\}}\mu_u
                      +\mathbf1_{\{w=v\}}\mu_v.
\]
Multiply by the nonnegative weight $\phi-w$. On the first contact set it equals $\phi-u$, and on the second it equals $\phi-v$. Approximation by bounded weights proves the displayed bound; polar sets have zero mass. Both right-hand terms are finite. The energy characterization proves $w\in\mathcal E^1$, and constant-shift invariance restores the original potentials.
\end{proof}
\index{rooftop envelope!finite-energy closure}

\begin{corollary}\label{cor:mutual-full-mass-singularity-envelopes}
For $u,v\in\mathcal E(X,\theta,\phi)$ one has $P_\theta[u](v)=v$.
\end{corollary}
\begin{proof}
Their rooftop $w=P_\theta(u,v)$ has relative full mass by the preceding mass argument. Since $w\leq v$ and $M(w)=M(v)>0$, \cref{prop:relative-envelope-full-mass-characterization} gives $P_\theta[w](v)=v$. Every competitor of singularity type at most $w$ has type at most $u$, so
$v=P_\theta[w](v)\leq P_\theta[u](v)\leq v$.
\end{proof}

\Needspace{16\baselineskip}
\begin{proposition}\label{prop:finite-energy-affine-first-variation}
For $u,v\in\mathcal E^1(X,\theta,\phi)$ and $u_t=(1-t)u+tv$, the energy is continuously differentiable on $[0,1]$, with one-sided derivatives at the endpoints, and
\[
 \frac{\mathrm d}{\mathrm dt}E_\theta^\phi(u_t)
                                  =\int_X(v-u)\,\mu_{u_t}.
\]
The derivative is a polynomial in $t$.
\end{proposition}
\begin{proof}
Cross-integrability makes every coefficient
$\int_X(v-u)\langle\theta_v^a\wedge\theta_u^{n-a}\rangle$ finite. Multilinearity makes the asserted derivative a polynomial. Apply the finite-energy cocycle to $u_{t+s},u_t$, divide by $s\ne0$, and use $u_{t+s}-u_t=s(v-u)$. Each mixed integral is a polynomial in $s$ with those finite coefficients. At $s=0$ each of the $n+1$ terms equals $\int_X(v-u)\mu_{u_t}$. This proves both derivatives and their continuity, including the endpoints.
\end{proof}

\begin{lemma}\label{lem:rooftop-affine-energy-derivative}
For $u,v\in\mathcal E^\infty(X,\theta,\phi)$ set $q_t=(1-t)u+tv$ and $\eta_t=P_\theta(q_t,v)$. Then
\[
 \frac{\mathrm d}{\mathrm dt}E_\theta^\phi(\eta_t)
                   =\int_X(v-u)_+\,\mu_{\eta_t},\qquad0\leq t\leq1,
\]
with one-sided endpoint derivatives. The energy is continuously differentiable.
\end{lemma}
\begin{proof}
The obstacle $f_t=\min\{q_t,v\}$ satisfies
$f_{t+s}-f_t=s(v-u)_+$ whenever both parameters lie in $[0,1]$, outside the common pole set. The obstacle calculus implies
$\|\eta_{t+s}-\eta_t\|_\infty\leq |s|\|u-v\|_\infty$ in the sense of bounded quasi-psh differences. All rooftops stay in one bounded-type interval about $\phi$.

For $s>0$, tangent bounds and contact concentration give
\[
 \begin{aligned}
 s\int_X(v-u)_+\,\mu_{\eta_{t+s}}
 &\leq E_\theta^\phi(\eta_{t+s})-E_\theta^\phi(\eta_t)\\
 &\leq s\int_X(v-u)_+\,\mu_{\eta_t}.
 \end{aligned}
\]
For example, on the lower-bound contact set $\eta_{t+s}=f_{t+s}$ and $\eta_t\leq f_t$; on the upper-bound contact set $\eta_t=f_t$ and $\eta_{t+s}\leq f_{t+s}$. Uniform difference convergence gives capacity convergence and hence weak measure convergence by the relative full-mass theorem. Strip rescaling and capacity comparison give uniform absolute continuity for these measures. The fixed bounded quasi-continuous weight $(v-u)_+$ therefore passes to the limit by the weighted convergence lemma. Divide the display by $s$ and let $s\downarrow0$. Interchanging $t$ and $t+s$ proves the left derivative. The same weighted convergence proves continuity of the derivative.
\end{proof}

\Needspace{16\baselineskip}
\begin{theorem}\label{thm:relative-energy-diamond}
For $u,v\in\mathcal E^1(X,\theta,\phi)$,
\begin{equation}\label{eq:relative-energy-diamond}
 E_\theta^\phi(\max\{u,v\})+E_\theta^\phi(P_\theta(u,v))
                      \geq E_\theta^\phi(u)+E_\theta^\phi(v).
\end{equation}
\end{theorem}
\begin{proof}
First let $u,v$ have bounded difference from $\phi$, and use the preceding $q_t,\eta_t$. Put $h=(v-u)_+$. For $t<1$, on $\{u<v\}$ one has $q_t<v$, so $\eta_t\leq q_t<v$. The rooftop contact inequality therefore gives
\[
 \int_Xh\,\mu_{\eta_t}
                   \leq\int_{\{u<v\}}(v-u)\,\mu_{q_t}.
\]
On the same plurifine open set, $(1-t)u+t\max\{u,v\}=q_t$. Affine variation and locality give
\[
 \begin{aligned}
 E_\theta^\phi(\max\{u,v\})-E_\theta^\phi(u)
  &=\int_0^1\int_{\{u<v\}}(v-u)\,\mu_{q_t}\,\mathrm dt\\
  &\geq E_\theta^\phi(v)-E_\theta^\phi(P_\theta(u,v)).
 \end{aligned}
\]
The possible endpoint $t=1$ does not affect the integral. This is the bounded-type inequality.

For finite-energy $u,v$, subtract a common constant to put them below $\phi$ and truncate both at $\phi-k$. Their maxima decrease to $\max\{u,v\}$, and their rooftops decrease to $P_\theta(u,v)$ by decreasing obstacle calculus. All four limits have finite energy by upper closure and rooftop closure. Decreasing energy continuity passes the bounded-type inequality to the limit. Restoring the common constant adds the same $2mC$ to both sides.
\end{proof}
\index{Monge--Amp\`ere energy!diamond inequality}

\begin{example}\label{ex:lone-convergence-energy-drop}
Energy need not be continuous for $L^1$ convergence, even for smooth potentials on $\mathbb P^1$. Let $\omega=\ddc\log(1+|z|^2)$, of mass one, and take $a\downarrow0$, $\varepsilon_a=\exp(-a^{-2})$ and
\[
 u_a(z)=a\log\frac{|z|^2+\varepsilon_a}{1+|z|^2}.
\]
These functions extend smoothly across infinity, and
$\omega+\ddc u_a=(1-a)\omega+a\ddc\log(|z|^2+\varepsilon_a)>0$ for $0<a<1$. They tend to zero in $L^1$: they lie between zero and the integrable function $a\log(|z|^2/(1+|z|^2))$.

Write $r=\log|z|^2$ and $F(r)=\log(e^r+\varepsilon)-\log(1+e^r)$. Radial integration with the normalization of \cref{thm:euclidean-jensen-lelong}, followed by Stokes and integration by parts in $r$, gives
\[
 \int_X\mathrm du_a\wedge\mathrm d^{\mathrm c}u_a
             =a^2\int_{-\infty}^\infty F'(r)^2\,\mathrm dr
             =a^2(-\log\varepsilon_a+O(1))\longrightarrow1.
\]
Here $0\leq F'\leq1$. On $[\log\varepsilon,0]$, its difference from $1$ is at most $\varepsilon e^{-r}+e^r$, whose integral is at most $2$. On the left tail $F'\leq e^r/\varepsilon$, and on the right tail $F'\leq e^{-r}$; the squared tail integrals are uniformly bounded. These estimates justify the uniform $O(1)$. The curve energy formula therefore gives $E_\omega(u_a)\to-1/2$, whereas $E_\omega(0)=0$.
\end{example}
